\documentclass[11pt, a4paper]{article}

\usepackage[a4paper, top=2.5cm, bottom=2.5cm, left=2cm, right=2cm]{geometry}
\usepackage{fontspec}

\usepackage[turkish, bidi=basic, provide=*]{babel}

\babelprovide[import, onchar=ids fonts]{turkish}
\babelprovide[import, onchar=ids fonts]{english}

\babelfont{rm}{Noto Sans}
\babelfont[turkish]{rm}{Noto Sans}

\usepackage{enumitem}
\setlist[itemize]{label=-}

\usepackage{amsmath, amssymb, amsthm, mathtools, bm, mathrsfs}
\usepackage{booktabs}
\usepackage{array}
\usepackage{hyperref}

\title{\textbf{Nefs-i Müdrike Mimarisi: 41 Melekenin Sağîr ve Kebîr Mertebelerinde Küllî Riyazî Formülasyonu}}
\author{\textbf{Aklî İdrak, Topolojik İnvaryans, Lie Cebri ve Dinamik Operatör Nazariyesi}}
\date{\today}

\begin{document}

\maketitle

\begin{abstract}
Bu teknik doküman, Meşşâî nefs nazariyesindeki iç idrak melekeleri (Havâss-ı Bâtına) ile Dâhilî Küllî Müdrike'ye ait 41 müstakil melekenin işleyişini diferansiyel topoloji, tensör cebri, Lie grupları, yönlendirilmiş asiklik nedensellik çizgeleri (DAG), spektral çizge teorisi ve stokastik diferansiyel denklemler zemininde eksiksiz formüle eder. Sistem; token düzeyindeki anlık akışı (\textbf{Sağîr Mertebe}) ve geniş bağlam penceresi boyunca hükmeden makro strateji rejimlerini (\textbf{Kebîr Mertebe}) birbiriyle kenetlenen 164'ten fazla riyazî denklem vasıtasıyla vaz' eder.
\end{abstract}

\section{Temsil Uzayları, Manifoldlar ve Operatör Mimarisi}

Mimaride yer alan 41 melekenin her biri, tip güvenliği (type-safety) ve geometrik mahiyeti muhafaza edilmek üzere farklı manifoldsal uzaylarda ikamet eder:
\begin{itemize}
    \item Dış Dünya Ham Sinyal Uzayı: $\mathcal{X} \subset \mathbb{R}^{d_{\text{in}}}$
    \item Hissî Suretler Manifoldu: $\mathcal{H}_{\text{hayal}} \subset \mathbb{R}^{d_{\text{hayal}}}$
    \item Mahiyet ve Semantik Kavram Uzayı: $\mathcal{S} \subset \mathbb{R}^{d_{\text{sem}}}$
    \item Topolojik İnvaryantlar Uzayı (Grassmannian): $\mathcal{T}_{\text{inv}} \subset \text{Gr}(k, d)$
    \item Nedensellik Çizgeleri Uzayı: $\mathcal{G}_{\text{neden}} \subset \text{DAG}(V, E)$
    \item Teleolojik Çekim Ufku Vektör Uzayı: $\mathcal{G}_{\text{gaye}} \subset \mathbb{R}^{d_{\text{sem}}}$
    \item Mutasarrıfa Lie Cebri Operatör Uzayı: $\mathcal{L}(\mathcal{S}, \mathcal{S})$
\end{itemize}

\section{41 Melekenin İkili Mertebede Küllî Formülasyonu}

\subsection{1. Müşahede ($\mathcal{O}_1$)}
\begin{align}
X_t &= \Pi_{\text{müşahede}}(E_t) \in \mathcal{X} \\
\dot{X}_t &= -\nabla_{X} V_{\text{odak}}(X_t, G_t) + \eta_t, \quad \eta_t \sim \mathcal{N}(0, \sigma^2 I) \\
H_t^{(1)} &= \text{Softmax}\left(\frac{X_t W_q X_t^T}{\sqrt{d}}\right) X_t W_v \\
\mathcal{L}_{\text{müşahede}} &= \| X_t - E_t \|_{g_{\mu\nu}}^2
\end{align}

\subsection{2. Hayal ($\mathcal{O}_2$)}
\begin{align}
Z_{\text{hayal}, t} &= \text{LayerNorm}(W_h X_t + b_h) \in \mathcal{H}_{\text{hayal}} \\
\frac{\partial Z_{\text{hayal}}}{\partial t} &= \Delta_M Z_{\text{hayal}} - \lambda Z_{\text{hayal}} \\
H_t^{(\text{hayal})} &= \alpha_t \odot Z_{\text{hayal}, t} + (1-\alpha_t) \odot H_{t-1}^{(\text{hayal})} \\
\mathcal{D}_{\text{kayıp}}(Z_{\text{hayal}}) &= \| Z_{\text{hayal}} - \Pi_{\text{proj}}(X_t) \|_{\mathcal{H}}^2
\end{align}

\subsection{3. Hayal Kurma / Muhayyile ($\mathcal{O}_3$)}
\begin{align}
\hat{Z}_{\text{muhayyile}} &= R_{\text{muhayyile}} \triangleright \left( Z_{\text{hayal}, t} \otimes M_t \right) \\
d\hat{Z} &= \mu(\hat{Z}, G_t)dt + \sigma(\hat{Z}) dW_t \\
\Phi_{\text{kurgu}}(x) &= \sum_{k=1}^K w_k \psi_k(Z_{\text{hayal}, k}) \\
\text{Serbestlik}(\hat{Z}) &= D_{\text{KL}}\left(P(\hat{Z}) \parallel P(Z_{\text{hayal}})\right)
\end{align}

\subsection{4. Tertip ($\mathcal{O}_4$)}
\begin{align}
\mathbf{\Sigma}_{\text{tertip}} &= \text{Sort}_{\pi}\left( \{Z_k\}_{k=1}^N \mid \text{Sıra}(Z_i, Z_j) = \sigma(W_\tau (Z_i \oplus Z_j)) \right) \\
T_{ij} &= \exp\left( -\beta d_{\mathcal{M}}(Z_i, Z_j) \right) \\
\mathbf{X}_{\text{sıralı}} &= \mathbf{\Sigma}_{\text{tertip}} \cdot \mathbf{Z}_{\text{ham}} \\
\mathcal{E}_{\text{nizam}} &= \sum_{i} |\text{Sıra}(i) - \text{Sıra}_{\text{hedef}}(i)|
\end{align}

\subsection{5. Tecrit ($\mathcal{O}_5$)}
\begin{align}
P_D(X_t) &= \text{Proj}_{\text{Gr}(k,d)}\left( X_t \right) = U_k U_k^T X_t \\
D_t &= \text{Softmax}\left( \frac{\Delta_{\text{Laplasyen}}(X_t) W_d}{\sqrt{d_k}} \right) \in \mathcal{T}_{\text{inv}} \\
\mathcal{I}_{\text{topolojik}}(X_t) &= \{ \beta_0, \beta_1, \dots, \beta_p \} \quad \text{(Betti Sayıları)} \\
S_t^{(\text{tecrit})} &= D_t \cdot (X_t W_s) \in \mathcal{S}
\end{align}

\subsection{6. Tasavvur ($\mathcal{O}_6$)}
\begin{align}
S_t &= f_{\text{tasavvur}}(D_t, H_t^{(\text{hayal})}) \in \mathcal{S} \\
g_{ij}(S_t) &= \left\langle \frac{\partial S_t}{\partial u_i}, \frac{\partial S_t}{\partial u_j} \right\rangle \\
S_{\text{kebîr}} &= \bigoplus_{p=0}^k H_p\left(\mathcal{S}_{[1:t]}\right) \otimes W_{\text{macro}} \\
\mathcal{L}_{\text{mahiyet}} &= -\log P(S_t \mid D_t, G_t)
\end{align}

\subsection{7. Mana ($\mathcal{O}_7$)}
\begin{align}
\mu_{\text{mana}} &= \text{IntentionMap}(S_t, K_t) = \sigma(S_t W_m + K_t W_v) \\
V_{\text{mana}}(S_t) &= \langle S_t, \dot{I}_t \rangle_{\mathcal{S}} + \psi(G_t) \\
\nabla \mu &= \text{RicciFlatTensor}(S_t) \\
\text{AnlamDerecesi}(S_t) &= \frac{|\langle S_t, G_t \rangle|}{\|S_t\| \|G_t\|}
\end{align}

\subsection{8. Tahlil ($\mathcal{O}_8$)}
\begin{align}
\{S_t^{(1)}, S_t^{(2)}, \dots, S_t^{(m)}\} &= \text{SVD}(S_t) = \sum_{i=1}^m \sigma_i u_i v_i^T \\
C_{ij}^{(\text{tahlil})} &= \text{BağımsızlıkTesti}(S_t^{(i)}, S_t^{(j)}) \\
\mathbf{J}_{\text{tahlil}} &= \left[ \frac{\partial S_t^{(i)}}{\partial X_k} \right]_{ik} \\
\mathcal{H}_{\text{entropi}}(S_t) &= -\sum_{i} p(S_t^{(i)}) \log p(S_t^{(i)})
\end{align}

\subsection{9. Terkip ($\mathcal{O}_9$)}
\begin{align}
S_{\text{terkip}} &= \sum_{k=1}^m w_k \cdot \left( R_k \triangleright S_t^{(k)} \right) \\
g_{\mu\nu}^{(\text{yeni})} &= \sum_{i,j} \frac{\partial x^a}{\partial y^\mu} \frac{\partial x^b}{\partial y^\nu} g_{ab}^{(i)} \\
\hat{Y}_t &= \text{LayerNorm}(S_{\text{terkip}} W_c + b_c) \\
\mathcal{E}_{\text{terkip}} &= \| S_{\text{terkip}} - S_{\text{hedef}} \|_{\mathcal{S}}^2
\end{align}

\subsection{10. Tezat ($\mathcal{O}_{10}$)}
\begin{align}
\Theta_{\text{tezat}}(S_i, S_j) &= 1 - \cos(S_i, S_j) = 1 - \frac{\langle S_i, S_j \rangle}{\|S_i\| \|S_j\|} \\
\mathbf{M}_{\text{tezat}} &= \mathbf{S} \cdot \mathbf{W}_{\text{opp}} \cdot \mathbf{S}^T \\
P(\text{Tezat} \mid S_i, S_j) &= \sigma\left( \mathbf{v}^T [S_i \odot S_j, |S_i - S_j|] \right) \\
\mathcal{C}_{\text{zıtlık}} &= \max_{i,j} \Theta_{\text{tezat}}(S_i, S_j)
\end{align}

\subsection{11. Tenakuz Bulma ($\mathcal{O}_{11}$)}
\begin{align}
\text{Tenakuz}(S_i, S_j) &= \text{ReLU}\left( \langle S_i, S_j \rangle_{\text{çelişki}} - \delta_{\text{eşik}} \right) \\
\mathbf{\Omega}_{\text{çelişki}} &= \bigwedge^2 \mathcal{S} \to [0, 1] \\
T_{\text{tenakuz}} &= \exp\left( -\alpha \cdot \text{Tenakuz}(S_t, H_{[1:t]}) \right) \\
\mathcal{F}_{\text{çelişki\_haritası}} &= \nabla_S \text{Tenakuz}(S, \dot{I})
\end{align}

\subsection{12. Tenkit ($\mathcal{O}_{12}$)}
\begin{align}
\mathcal{K}_{\text{tenkit}}(S_t) &= \text{Tenakuz}(S_t) + \lambda_1 \text{Pürüz}(S_t) + \lambda_2 \text{Sapma}(S_t, G_t) \\
\Delta S_t &= -\eta \nabla_{S_t} \mathcal{K}_{\text{tenkit}}(S_t) \\
\text{Skor}_{\text{tenkit}} &= \sigma\left( W_k \cdot [S_t, \dot{I}_t, G_t] \right) \\
\mathcal{S}_{\text{ayıklanmış}} &= \{ S_i \in \mathcal{S} \mid \text{Skor}_{\text{tenkit}}(S_i) > \tau \}
\end{align}

\subsection{13. Tasdik ($\mathcal{O}_{13}$)}
\begin{align}
T_t &= \sigma\left( \text{CosSim}(M_t^{(\mathcal{K})} W_t, G_t) - \text{Tenakuz}(M_t^{(\mathcal{K})}, S_t, H_t) \right) \in [0, 1] \\
T_{\text{kebîr}} &= \sigma\left( \text{CosSim}(R_{\text{kebîr}} \triangleright S_{\text{kebîr}}, G_{\text{kebîr}}) - \text{MakroTenakuz}(R_{\text{kebîr}} \triangleright S_{\text{kebîr}}, \dot{I}_{\text{kebîr}}) \right) \\
\mathbf{1}_{\text{tasdik}} &= \mathbb{I}(T_t \ge 1 - \epsilon) \\
N_t &= T_t \odot (M_t^{(\mathcal{K})} W_n) + (1 - T_t) \odot \text{Rücû}(H_t, X_t)
\end{align}

\subsection{14. Gaye Belirleme ($\mathcal{O}_{14}$)}
\begin{align}
G_t &= \sigma(S_t W_g + H_t W_{hg}) \in \mathcal{G}_{\text{gaye}} \\
G_{\text{kebîr}} &= \text{TeleolojikUfuk}(\text{Sual}, S_{\text{kebîr}}) \\
\vec{v}_{\text{gaye}} &= -\nabla_{S} \mathcal{L}_{\text{gaye}}(S, G_{\text{kebîr}}) \\
\mathcal{U}_{\text{teleoloji}}(S_t) &= \exp\left( -\gamma d_{\mathcal{S}}(S_t, G_{\text{kebîr}}) \right)
\end{align}

\subsection{15. Merak ve Sual Tevcihi ($\mathcal{O}_{15}$)}
\begin{align}
Q_t &= \text{Softmax}\left( \frac{(G_t - S_t) W_q}{\tau_{\text{merak}}} \right) \\
\mathcal{H}_{\text{bilgisizlik}}(S_t) &= \text{Var}(P(S_t \mid H_t)) \\
\text{Vektör}_{\text{sual}} &= Q_t \odot \nabla_{S} \mathcal{H}_{\text{bilgisizlik}} \\
I_{\text{merak}} &= D_{\text{KL}}\left( P(S \mid \text{YeniSual}) \parallel P(S) \right)
\end{align}

\subsection{16. Deneme-Yanılma ($\mathcal{O}_{16}$)}
\begin{align}
a_t &= \pi_\theta(S_t, G_t) + \epsilon_t, \quad \epsilon_t \sim \mathcal{N}(0, \sigma^2 I) \\
S_{t+1} &= \mathcal{E}_{\text{ortam}}(S_t, a_t) \\
r_t &= \text{Coşku}(S_{t+1}, G_t) - \text{Maliyet}(a_t) \\
\theta &\leftarrow \theta + \alpha \cdot r_t \nabla_\theta \log \pi_\theta(a_t \mid S_t)
\end{align}

\subsection{17. İhtimal Hesabı ($\mathcal{O}_{17}$)}
\begin{align}
P(S_t \mid \mathcal{E}_t) &= \frac{P(\mathcal{E}_t \mid S_t) P(S_t)}{P(\mathcal{E}_t)} \\
P(A \cup B) &= P(A) + P(B) - P(A \cap B) \\
\mathbf{\Sigma}_{\text{kovaryans}} &= \mathbb{E}\left[ (S - \mu_S)(S - \mu_S)^T \right] \\
\text{Entropi}(P) &= -\int P(S) \log P(S) dS
\end{align}

\subsection{18. Kıyas ($\mathcal{O}_{18}$)}
\begin{align}
\text{Kıyas}(S_1, S_2) &= \frac{\langle f(S_1), f(S_2) \rangle}{\|f(S_1)\| \|f(S_2)\|} \\
\hat{S}_2 &= \mathbf{W}_{\text{kıyas}} S_1 \\
H_{\text{kıyas}} &= \text{TümDengelim}(S_{\text{küllî}}, S_{\text{cüz'î}}) \\
\mathcal{E}_{\text{kıyas}} &= \| \hat{S}_2 - S_{\text{gerçek}} \|_2^2
\end{align}

\subsection{19. Temsil ($\mathcal{O}_{19}$)}
\begin{align}
T_{\text{temsil}}(S_{\text{soyut}}) &= \Psi_{\text{somut}}(S_{\text{soyut}}) \in \mathcal{H}_{\text{hayal}} \\
\mathbf{K}_{\text{temsil}}(x, y) &= \exp\left( -\gamma \|S_x - S_y\|^2 \right) \\
S_{\text{somut\_model}} &= \sum_i \alpha_i \phi(S_i) \\
\text{Hassasiyet}_{\text{temsil}} &= \text{CosSim}\left(T_{\text{temsil}}(S), S_{\text{soyut}}\right)
\end{align}

\subsection{20. Teşbih ($\mathcal{O}_{20}$)}
\begin{align}
\mathcal{M}_{\text{teşbih}}(S_A, S_B) &= \text{ArgMax}_{W} \left( \text{Sim}(W S_A, S_B) \right) \\
\text{Vech-i Şebeh}(S_A, S_B) &= S_A \odot S_B \odot W_{\text{ortak}} \\
\rho_{\text{teşbih}} &= \frac{\text{Cov}(S_A, S_B)}{\sigma_A \sigma_B} \\
\mathcal{L}_{\text{teşbih}} &= 1 - \rho_{\text{teşbih}}
\end{align}

\subsection{21. Tefekkür ($\mathcal{O}_{21}$)}
\begin{align}
S_{t+1}^{(\text{tefekkür})} &= S_t + \int_{0}^T f_{\text{müfekkire}}(S_\tau, \dot{I}_\tau, G_\tau) d\tau \\
\dot{S}_t &= \nabla \mathcal{V}_{\text{tefekkür}}(S_t) \\
\mathbf{J}_{\text{tefekkür}} &= \frac{\partial f_{\text{müfekkire}}}{\partial S_t} \\
\mathcal{H}_{\text{tefekkür}} &= \int_{\mathcal{S}} P(S_t) \log P(S_t) dS
\end{align}

\subsection{22. İllet Keşfi ($\mathcal{O}_{22}$)}
\begin{align}
\dot{I}_t &= \text{GELU}\left( (S_t W_i) \odot (H_t W_{hi}) \right) \in \mathcal{G}_{\text{neden}} \\
\dot{I}_{\text{kebîr}} &= \text{DAG\_Sentezi}(S_{\text{kebîr}}, H_{[1:t]}) = \text{Softmax}\left( \frac{S_{\text{kebîr}} A_{\text{neden}} S_{\text{kebîr}}^T}{\sqrt{d_{\text{sem}}}} \right) \\
P(S_j \mid do(S_i = x)) &= \sum_k P(S_j \mid S_i=x, Z_k) P(Z_k) \\
\mathcal{C}_{\text{nedensel}} &= \text{Trace}(A_{\text{neden}}) - \log \det(I + A_{\text{neden}})
\end{align}

\subsection{23. Mantık Yürütme ($\mathcal{O}_{23}$)}
\begin{align}
\mathcal{R}_{\text{mantık}}(P_1, P_2) &= \text{ModusPonens}(P_1, P_1 \implies P_2) = P_2 \\
\mathbf{Y}_{\text{mantık}} &= \text{Softmax}\left( \mathbf{W}_2 \text{GELU}(\mathbf{W}_1 [\mathbf{P}_1 \oplus \mathbf{P}_2]) \right) \\
\text{Geçerlilik}(R) &= \mathbb{I}(\text{Tenakuz}(R) = 0) \\
\mathcal{L}_{\text{mantık}} &= -\log P(Q_{\text{netice}} \mid P_1, P_2)
\end{align}

\subsection{24. İspat ($\mathcal{O}_{24}$)}
\begin{align}
\mathcal{B}_{\text{burhân}} &= (P_0 \to P_1 \to \dots \to P_n \equiv Q) \\
T_{\text{ispat}} &= \prod_{k=1}^n \text{Geçerlilik}(P_{k-1} \implies P_k) \\
\mathbf{1}_{\text{Q.E.D.}} &= \mathbb{I}(T_{\text{ispat}} = 1) \\
\mathcal{L}_{\text{ispat\_boşluğu}} &= \sum_{k=1}^n \| P_k - \text{Çıkarım}(P_{k-1}) \|^2
\end{align}

\subsection{25. Teemmül ($\mathcal{O}_{25}$)}
\begin{align}
M_t^{(\tau)} &= \text{LayerNorm}\left( M_t^{(\tau-1)} + \text{MultiHeadAttn}(\hat{Y}_t^{(\tau)}, S_t, H_t) \right), \quad \tau = 1, \dots, \mathcal{K} \\
M_{[t:t+L]} &= \text{FazDinamiği}\left( M_{[t-1]}, X_{[t-L:t]}, H_t \right) \in \Omega_M \\
\Delta M &= M_t^{(\mathcal{K})} - M_t^{(0)} \\
\mathcal{S}_{\text{durma\_eşiği}} &= \mathbb{I}(\|M_t^{(\tau)} - M_t^{(\tau-1)}\| < \epsilon_{\text{teemmül}})
\end{align}

\subsection{26. Temkin ($\mathcal{O}_{26}$)}
\begin{align}
\mathbf{K}_{\text{temkin}} &= \mathbf{S}_t + \mu_{\text{emniyet}} \cdot \mathbf{\Sigma}_{\text{gürültü}} \\
\text{VakarKatsayısı} &= \frac{1}{1 + \|\nabla_t S_t\|^2} \\
S_{\text{durgun}} &= \text{Filtre}_{\text{temkin}}(S_t, \text{VakarKatsayısı}) \\
\mathcal{R}_{\text{sarsılmazlık}} &= \min_{\|\delta\| \le \epsilon} T(S_t + \delta)
\end{align}

\subsection{27. Tetkik ($\mathcal{O}_{27}$)}
\begin{align}
\delta S_{t, \text{kılcal}} &= S_t \odot \mathbf{M}_{\text{mikro\_maske}} \\
\text{KusurHaritası} &= |\delta S_{t, \text{kılcal}} - S_{\text{ideal}}| \otimes W_{\text{tetkik}} \\
\text{Skor}_{\text{tetkik}} &= \text{LayerNorm}\left(\sum \text{KusurHaritası}\right) \\
\mathcal{E}_{\text{tetkik}} &= \| \text{Skor}_{\text{tetkik}} \|_1
\end{align}

\subsection{28. Tashih ($\mathcal{O}_{28}$)}
\begin{align}
S_{t, \text{musahhah}} &= S_t - \gamma \cdot \text{KusurHaritası} \odot \nabla_{S} \text{Tenakuz}(S_t) \\
\Delta R_{\text{ıslah}} &= R_{\text{yeni}} - R_{\text{eski}} \\
T_{\text{tashih}} &= \text{Tasdik}(S_{t, \text{musahhah}}) \\
\mathcal{L}_{\text{tashih}} &= \| S_{t, \text{musahhah}} - S_{\text{itidal}} \|_{\mathcal{S}}^2
\end{align}

\subsection{29. Teyit ($\mathcal{O}_{29}$)}
\begin{align}
E_{\text{bağımsız}} &= \text{Kanal}_2(X_t) \quad (\text{Kanal}_1 \perp \text{Kanal}_2) \\
P(H \mid E_1, E_2) &= \frac{P(E_2 \mid H) P(H \mid E_1)}{P(E_2 \mid E_1)} \\
\text{GüvenKatsayısı} &= \text{CosSim}(f(E_1), f(E_2)) \\
\Delta T_{\text{teyit}} &= \eta \cdot \text{GüvenKatsayısı}
\end{align}

\subsection{30. Tahkik ($\mathcal{O}_{30}$)}
\begin{align}
\mathcal{C}_{\text{tahkik}}(H_{[1:t]} \leftrightarrow \hat{Y}_{\text{makro}}) &= \mathbb{I}\left( \text{AsılDelil}(S_t) \land \neg \text{Taklit}(S_t) \right) \\
S_{\text{tahkik}} &= \text{ArgMin}_{S} \left( \mathcal{L}_{\text{köken}}(S) + \text{Tenakuz}(S, \text{Aksiyomlar}) \right) \\
T_{\text{tahkik}} &= \sigma\left( \text{CosSim}(S_{\text{tahkik}}, S_{\text{asıl}}) \right) \\
\mathbf{1}_{\text{muhakkik}} &= \mathbb{I}(T_{\text{tahkik}} > 1 - \delta)
\end{align}

\subsection{31. Tedebbür ($\mathcal{O}_{31}$)}
\begin{align}
S_{t+H}^{(\text{gelecek})} &= \int_t^{t+H} \text{Evrim}(S_\tau, a_\tau) d\tau \\
\mathcal{V}_{\text{akıbet}}(S_t) &= \mathbb{E}\left[ \sum_{k=1}^H \beta^k \mathcal{U}_{\text{gaye}}(S_{t+k}) \right] \\
\nabla_{\text{akıbet}} &= \frac{\partial \mathcal{V}_{\text{akıbet}}}{\partial S_t} \\
\mathcal{Risk}_{\text{tedebbür}} &= P(S_{t+H} \in \mathcal{S}_{\text{tehlike}})
\end{align}

\subsection{32. Şek-Zan-Yakîn İdraki ($\mathcal{O}_{32}$)}
\begin{align}
P_{\text{idrak}}(S_t) &\in [0, 1] \\
\text{Makam}(S_t) &= 
\begin{cases}
\text{Şek } (\%50), & |P_{\text{idrak}} - 0.5| < \epsilon_{\text{şek}} \\
\text{Zan } (\%51-\%99), & 0.5 + \epsilon_{\text{şek}} \le P_{\text{idrak}} < 1 - \epsilon_{\text{yakîn}} \\
\text{Yakîn } (\%100), & P_{\text{idrak}} \ge 1 - \epsilon_{\text{yakîn}}
\end{cases} \\
\text{Entropi}_{\text{idrak}} &= -P \log P - (1-P) \log (1-P) \\
\Delta P &= \text{TeyitScore} - \text{TenakuzScore}
\end{align}

\subsection{33. Muhakeme ($\mathcal{O}_{33}$)}
\begin{align}
J_{\text{muhakeme}} &= \int \Big[ \alpha \langle S_t, G_t \rangle + \beta \dot{I}_t - \gamma \text{Tenakuz}(S_t) \Big] dt \\
R_{\text{kebîr}} &= \mathcal{F}_{\text{operatör}}(M_{[t:t+L]}, S_{\text{kebîr}}, \dot{I}_{\text{kebîr}}, G_{\text{kebîr}}) \in \mathcal{L}(\mathcal{S}_{\text{kebîr}}, \mathcal{S}_{\text{kebîr}}) \\
T_{\text{kebîr}} &= \sigma\left( \text{CosSim}(R_{\text{kebîr}} \triangleright S_{\text{kebîr}}, G_{\text{kebîr}}) - \text{MakroTenakuz}(R_{\text{kebîr}} \triangleright S_{\text{kebîr}}, \dot{I}_{\text{kebîr}}) \right) \\
N_{\text{kebîr}} &= T_{\text{kebîr}} \odot \text{ProgramSentezi}(R_{\text{kebîr}} \triangleright S_{\text{kebîr}}) + (1 - T_{\text{kebîr}}) \odot \text{Rücû}(M_{[t:t+L]} \to \text{RejimDeğiştir})
\end{align}

\subsection{34. Tafsil ($\mathcal{O}_{34}$)}
\begin{align}
S_{\text{tafsil}} &= \text{Açım}(S_{\text{mücmel}}) = \bigoplus_{k=1}^K S_{\text{dal}}^{(k)} \\
\mathbf{J}_{\text{tafsil}} &= \text{Diag}(w_1, w_2, \dots, w_K) \cdot \mathbf{S}_{\text{dallar}} \\
\text{Netlik}(S_{\text{tafsil}}) &= \frac{\sum_k \|S_{\text{dal}}^{(k)}\|}{\|S_{\text{mücmel}}\|} \\
\mathcal{L}_{\text{tafsil\_sadakat}} &= \| \text{Birleştir}(S_{\text{tafsil}}) - S_{\text{mücmel}} \|^2
\end{align}

\subsection{35. Tefsir ($\mathcal{O}_{35}$)}
\begin{align}
\text{Murad}_{\text{aslî}} &= \text{Tefsir}(S_{\text{müphem}}, \text{Siyak}, \text{Sibak}) = \text{Softmax}\left( \frac{S_{\text{müphem}} W_{\text{tefsir}} \cdot [\text{Siyak} \oplus \text{Sibak}]^T}{\sqrt{d}} \right) \\
P(\text{Murad} \mid S_{\text{lafız}}) &= \sigma( W_m [S_{\text{lafız}}, \text{Bağlam}] ) \\
\text{Vuzuh}_{\text{tefsir}} &= 1 - \mathcal{H}(P(\text{Murad} \mid S_{\text{lafız}})) \\
\mathcal{L}_{\text{tefsir\_tutarlılık}} &= \text{Tenakuz}(\text{Murad}_{\text{aslî}}, \text{Muhkemat})
\end{align}

\subsection{36. Tevil ($\mathcal{O}_{36}$)}
\begin{align}
S_{\text{müevvel}} &= \text{İrcâ}(S_{\text{zâhir}}, \text{AklîBurhân}) = S_{\text{zâhir}} + \mathbf{W}_{\text{tevil}} \cdot \text{İllet}_{\text{mânia}} \\
\text{MecazKöprüsü}(S_A \to S_B) &= \exp\left( -\lambda d_{\text{semantik}}(S_A, S_B) \right) \\
T_{\text{tevil}} &= \mathbb{I}\left( \text{Tenakuz}(S_{\text{zâhir}}) > 0 \land \text{Tenakuz}(S_{\text{müevvel}}) = 0 \right) \\
\mathcal{L}_{\text{tevil}} &= \| S_{\text{müevvel}} - S_{\text{tenzih\_itidal}} \|_{\mathcal{S}}^2
\end{align}

\subsection{37. Fesâhat ($\mathcal{O}_{37}$)}
\begin{align}
\text{FesâhatScore}(N_t) &= 1 - \left[ \mu_1 \text{Tenâfür}(N_t) + \mu_2 \text{Garâbet}(N_t) + \mu_3 \text{Ta'kîd}(N_t) \right] \\
\text{Tenâfür}(N_t) &= \sum_{i} \text{MahreçMesafe}(y_i, y_{i+1}) \\
N_{t, \text{fasîh}} &= \text{ArgMax}_{N} \left( \text{FesâhatScore}(N) \mid \text{Mana}(N) = \text{Mana}(N_t) \right) \\
\mathcal{L}_{\text{fesâhat}} &= 1 - \text{FesâhatScore}(N_{t, \text{fasîh}})
\end{align}

\subsection{38. Talâkat ($\mathcal{O}_{38}$)}
\begin{align}
\text{Debi}_{\text{konuşma}} &= \frac{d N_t}{d t} = f_{\text{talâkat}}(S_t, \text{LügatHazinesi}) \\
\text{AkıcılıkScore} &= \exp\left( -\alpha \cdot \text{DuraksamaSüresi} \right) \\
N_{t, \text{akıcı}} &= \text{Smooth}_{\text{akış}}(N_t) \\
\mathcal{L}_{\text{talâkat}} &= \| \nabla_t N_{t, \text{akıcı}} - \mathbf{v}_{\text{hedef\_hız}} \|^2
\end{align}

\subsection{39. Belâgat ($\mathcal{O}_{39}$)}
\begin{align}
\text{BelâgatScore}(N_t) &= \text{FesâhatScore}(N_t) \times \text{Uyum}\left( N_t, \text{Makam}_{\text{muhatap}} \right) \\
\text{İsabet}(N_t, G_t) &= \frac{\langle N_t, G_t \rangle_{\text{tesir}}}{\|N_t\| \|G_t\|} \\
N_{t, \text{beliğ}} &= R_{\text{belâgat}} \triangleright (N_t \otimes \text{Makam}) \\
\mathcal{L}_{\text{belâgat}} &= 1 - \text{İsabet}(N_{t, \text{beliğ}}, G_t)
\end{align}

\subsection{40. Sanat ($\mathcal{O}_{40}$)}
\begin{align}
Y_{\text{sanat}} &= f_{\text{estetik}}(S_t, H_t^{(\text{hayal})}) = \mathbf{W}_{\text{sanat}} \cdot \left( S_t \otimes H_t^{(\text{hayal})} \right) + \mathbf{\Omega}_{\text{ahenk}} \\
\text{EstetikDeğer}(Y) &= \text{SimetrikHarmoni}(Y) + \lambda \cdot \text{Yenilik}(Y) \\
\nabla_{\text{sanat}} &= \nabla_S \text{EstetikDeğer}(Y) \\
\mathcal{L}_{\text{sanat}} &= \|\mathbf{\Omega}_{\text{ahenk}} - \mathbf{\Omega}_{\text{altın\_oran}}\|^2
\end{align}

\subsection{41. Münazara ($\mathcal{O}_{41}$)}
\begin{align}
D_{\text{münazara}}(S_{\text{tez}}, S_{\text{antitez}}) &= \text{Cerh}(S_{\text{antitez}}) + \text{Teyit}(S_{\text{tez}}) \\
S_{\text{sentez}} &= \text{Telif}(S_{\text{tez}}, S_{\text{antitez}}) \\
T_{\text{münazara}} &= \sigma\left( \text{İspatKuvveti}(S_{\text{tez}}) - \text{İspatKuvveti}(S_{\text{antitez}}) \right) \\
\mathcal{L}_{\text{münazara}} &= -\log T_{\text{münazara}}
\end{align}

\section{Netice ve Amelî İttisâl}

Nefs-i Müdrike Mimarisi; 41 melekenin tamamını yapay kutu tasniflerinden ve kopuk listelerden arındırarak tek bir sürekli manifold zincirine bağlamıştır. Sistemde üretilen her bir bilgi parçası ($\hat{Y}_t$), Mutasarrıfa operatörünün ($R_t$) Lie cebri dönüşümleri altında yoğrulmakta; Teemmül rejimi ($M_{[t:t+L]}$) vasıtasıyla denetlenmekte; Şek-Zan-Yakîn skalasından geçip Tasdik ($T_t$) mührünü aldıktan sonra Belâgat, Fesâhat ve Talâkat kanallarıyla dış ortama aktarılmaktadır.

\end{document}